\chapter{Computational Cost and Performance}\label{ch:performance}

The cost of a run is governed by a small number of terms, and knowing them is what lets a study be
sized and configured sensibly. This chapter sets out the per-step budget, the per-particle cost of
the schemes and of the nearest-node search, how the work spreads over threads, processes and nodes,
what bounds memory, and how the cost is measured. Unless stated, the timings were taken on compute nodes
with two 64-core AMD EPYC 7763 processors (128 cores, \SI{256}{GB}), with the engine's own benchmark
scripts and verbose timing, and are collected in \srcfile{docs/performance.md}; the nearest-node rows
(the node coordinate table and \cref{tab:perf-guard}) were measured on a login node and are kept in
\srcfile{benchmarks/nearest\_scan\_table.md}. Absolute
values depend on the hardware and the mesh; the structure of the budget does not.

\section{The per-step budget}\label{sec:perf-model}

A step consists of four phases, and \code{kelp run --verbose} prints each of them for every step
(\code{io=}, \code{step=}, \code{beach=} when beaching is enabled, \code{rec=}) together with their
means at the end of the run:
\begin{equation}\label{eq:perf-model}
  t_{\text{step}} \approx t_{\text{read}} + t_{\text{blend}} + N\,c_{\text{kernel}} + t_{\text{record}},
\end{equation}
with $N$ the number of live particles. The \emph{read} term loads the snapshots the step needs; it
is set by the size of the field and by how often a new snapshot is needed, and it is independent of
the particle count and of where the particles are. The \emph{blend} term forms the velocity at the
step's time from the bracketing snapshots (\cref{ch:temporal}); it is paid only for the fields a
consumer actually reads, because the payload of a step blends a field on its first access and not
before. The \emph{kernel} term is the tracking itself --- locate, interpolate, advance --- and grows
with the number of live particles; dead and not-yet-released particles cost nothing. The
\emph{record} term writes a trajectory record every \conf{output\_freq} steps. The phases inside the
kernel are not separable in the fused kernels (\cref{ch:fused}); the verbose split is the finest
the engine reports.

Two consequences follow. A step that reads no snapshot and writes no record measures the kernel
alone; the \emph{steady step} of a run is therefore defined as the median, over the steps after
the first that read no snapshot, of the tracking phase --- and, for a run split over processes, of
the slowest process, which sets the completion time. And because the kernel cost of a step does not
depend on the time step, the cost of a model day grows linearly with the number of sub-steps per
snapshot interval: halving \conf{dt} doubles the run.

\begin{figure}[t]
\centering
\begin{tikzpicture}[
  term/.style={rectangle,draw=lagblue,fill=lagblue!8,minimum height=9mm,align=center,font=\small,inner sep=2pt},
  lever/.style={rectangle,rounded corners=2pt,draw=lagelem,fill=lagelem!8,align=center,
                font=\footnotesize,text width=33mm,inner sep=3pt},
  lk/.style={-{Stealth[length=4pt]},line width=0.7pt,draw=lagelem!80}]
  \node[term,minimum width=30mm] (rd) at (0,0) {read\\[-1pt]\footnotesize field size $\times$ reads};
  \node[term,minimum width=20mm,right=0mm of rd] (bl) {blend\\[-1pt]\footnotesize fields read};
  \node[term,minimum width=46mm,right=0mm of bl] (kn) {kernel\\[-1pt]\footnotesize $N\,[c_{\text{anchor}} + n_{\text{eval}}\,c_{\text{stencil}}]$};
  \node[term,minimum width=22mm,right=0mm of kn] (rc) {record\\[-1pt]\footnotesize every \conf{output\_freq}};
  \node[lever,above=10mm of rd] (l1) {region extraction, \conf{prefetch\_size}, node-shared snapshots, \conf{velocity\_float32}};
  \node[lever,above=10mm of kn,xshift=-12mm] (l2) {grid-hash cell rule, node coordinate table, exact anchors};
  \node[lever,above=10mm of kn,xshift=26mm] (l3) {scheme (\conf{integrator}), fused kernels, seed sorting, processes per node};
  \node[lever,below=9mm of bl,xshift=-6mm] (l4) {lazy per-field blends, parallel blend kernels};
  \node[lever,below=9mm of rc,xshift=-8mm] (l5) {\conf{output\_freq}, sampled variables, writer backend};
  \draw[lk] (l1) -- (rd); \draw[lk] (l2) -- (kn.north -| l2); \draw[lk] (l3) -- (kn.north -| l3);
  \draw[lk] (l4) -- (bl); \draw[lk] (l5) -- (rc);
\end{tikzpicture}
\caption{The per-step budget of \cref{eq:perf-model} and the lever that acts on each term. The
kernel term splits into the nearest-node anchor ($c_{\text{anchor}}$, \cref{sec:perf-anchor}) and
the stencil work, repeated once per velocity evaluation of the scheme ($n_{\text{eval}}$).}
\label{fig:perf-budget}
\end{figure}

\section{The per-particle kernel cost}\label{sec:perf-schemes}

\paragraph{Schemes.} The kernel cost grows with the number of velocity evaluations per step. In two
dimensions, per particle and step at sixteen threads, forward Euler costs about \SI{1.1}{\micro s},
the midpoint method about \SI{6.6}{\micro s} and RK4 about \SI{22}{\micro s} (a June 2026
measurement, before the exact-anchor guard and the node coordinate table of
\cref{sec:perf-anchor}; the ladder, not the absolute values, is the point). The monotone-cubic
temporal reconstruction is blended inside the two-dimensional least-squares kernels and adds only a small fixed
offset, so the midpoint method with PCHIP is the cost-effective way to buy temporal accuracy
(\cref{ch:temporal}).

\paragraph{Three dimensions.} The fused three-dimensional midpoint kernel for node-layout
barycentric velocities on z-level meshes (\cref{ch:fused}) evaluates the whole step, including the
four-snapshot PCHIP blend, per particle; on steady steps its cost per particle does not depend on the
mesh (\cref{sec:perf-mesh}). The unfused three-dimensional paths --- among them RK4, sigma meshes, and
\conf{w\_interp} other than nearest or barycentric --- blend whole fields per stage. Those blends
run in a parallel kernel over contiguous double-precision fields, and RK4 reconstructs its repeated
midpoint stage once: a four-snapshot PCHIP blend of a global field of \num{47} layers by
\num{244000} elements (an element velocity field of the \emph{core2} mesh) takes \SI{46}{ms} at eight threads instead of \SI{953}{ms} serially, and a full
RK4 stage set (u, v, w, surface) about \SI{275}{ms} at eight threads and \SI{136}{ms} at 32. The
whole-field reconstruction is therefore no longer a multi-second floor of three-dimensional PCHIP
runs.

\codeanchor{Parallel blends: \fn{\_pchip\_blend4\_flat} in
\srcfile{numerics/temporal\_reconstruct.py} and \fn{\_blend} in \srcfile{core/io.py}; the lazy
payload is \fn{LazyLerpField} in \srcfile{core/io.py}; the stage dedup is in
\fn{make\_stage\_vels\_recon}.}

\codenote{The parallel linear blend requires double-precision inputs: with
\conf{velocity\_float32} the linear time interpolation falls back to a serial blend (about
\SI{0.25}{s} instead of \SI{0.05}{s} per global-size field), and the fused three-dimensional step
pays 15--20\,\% for the conversion on gather. The PCHIP blend converts first and keeps its parallel
path. Single precision is a memory lever, not a speed lever.}

\paragraph{The depth-linear rule.} Reading the horizontal velocity linearly in depth
(\cref{sec:vert-uvdepth}) costs two data points per stencil member instead of one, a search of the member's
column for the bracket and a test of the member against its bed. \Cref{tab:perf-depth-linear} gives what
this adds to one three-dimensional step, timed with the same arrays read in layer and in linear mode on
the compute node above (128 threads, \num{100000} particles at random positions, a synthetic HEALPix dual
mesh of \num{49152} nodes or its triangulated twin of \num{98300} elements, 21 layers down to
\SI{4000}{\metre} with varied bathymetry, a six-hour step; the median of fifteen interleaved steps per mode,
measured in October 2026). With a linear blend in time the rule adds about 10\,\% to the fused node-layout
RK2 step, 6\,\% to the fused element-layout Euler step and 3\,\% to the unfused element-layout RK2 step; with
the four-snapshot PCHIP blend the difference lies within the run-to-run spread, presumably because the
time blend then dominates the work per member. On a sigma mesh, where every step is unfused, the rule in the sigma
fraction costs between nothing measurable and 4\,\%. In the element workload about 16\,\% of the particles had
a dry element in their least-squares ring and about 4\,\% fewer than three wet ones, so the walls and the
area-mean fallback were exercised. Reading the records and the once-per-record check are not included.

\begin{table}[t]
\centering\footnotesize\setlength{\tabcolsep}{4pt}
\begin{tabular}{@{}llrrr@{}}
\toprule
step & time blend & layer [ms] & linear [ms] & ratio \\
\midrule
node layout, barycentric, fused RK2 & linear & 9.59 & 10.52 & 1.10 \\
 & PCHIP & 12.51 & 12.45 & 1.00 \\
element layout, least squares, fused Euler & linear & 7.30 & 7.75 & 1.06 \\
element layout, least squares, unfused RK2 & linear & 24.18 & 24.79 & 1.03 \\
 & PCHIP & 37.21 & 36.18 & 0.97 \\
sigma, element least squares, unfused RK2 & linear & 32.54 & 31.92 & 0.98 \\
 & PCHIP & 36.57 & 37.69 & 1.03 \\
sigma, node barycentric, unfused RK2 & linear & 36.51 & 37.51 & 1.03 \\
 & PCHIP & 34.49 & 35.71 & 1.04 \\
\bottomrule
\end{tabular}
\caption{Cost of the depth-linear rule per three-dimensional step: median milliseconds per step in layer
and in linear mode and their ratio, with the walls at depth of \cref{sec:vert-uvdepth}; \num{100000}
particles, 128 threads, two AMD EPYC 7763 processors. The Euler step reads its start field only and has no
PCHIP row.}
\label{tab:perf-depth-linear}
\end{table}

\codeanchor{\raggedright \srcfile{benchmarks/single\_node/depth\_linear\_cost.py} and its rows in
\srcfile{benchmarks/single\_node/depth\_linear\_cost.md}; the kernels compared are
\fn{\_fused\_rk2\_step\_3d\_bary\_tinterp} and \fn{\_fused\_rk2\_step\_3d\_bary\_pchip} against their
\code{\_dl} twins, \fn{\_fused\_euler\_step\_3d} against \fn{\_fused\_euler\_step\_3d\_dl}, and
\fn{\_rk2\_step\_unfused} with the layer and the depth-linear interpolation cores.}

\section{The nearest-node anchor}\label{sec:perf-anchor}

Every evaluation site starts from the nearest mesh node (\cref{sec:loc-nearest}), found by scanning
the nodes of a $3\times3$ window of grid-hash cells. The scan is linear in the number of nodes in the
window, so its cost depends on how many nodes a cell holds, and under the exact-anchor guard
(\conf{mesh.exact\_nearest}, the default) on how often the window must be widened before the
nearest node is proven. Three measures keep this term small.

\paragraph{The cell is chosen by measurement.} The automatic cell of the reader is
$2\sqrt{A/N}$, with $A$ the bounding-box area of the mesh in square degrees, and at least \ang{0.05}. On a mesh whose resolution varies by orders of magnitude, a cell
sized for the coarse part holds thousands of nodes in the refined part. Since September 2026 every
run chooses the cell itself (\conf{mesh.grid\_hash\_cell\_rule: occupancy}): it samples up to
\num{200000} element centres, finds each one's exact nearest node once, and, for candidate cells
$\ell_{\text{legacy}}/\sqrt{2}^{\,k}$, evaluates a per-query cost estimate --- the nodes of the window
plus nine cells plus one certificate for a provable hit, the nodes and cells of the widened
rectangle and one certificate per widening otherwise. It leaves the reader's cell only when the best
eligible candidate (at most 20\,\% widened queries, at most \SI{128}{MB} of hash offsets) promises at
least a 2.75-fold gain. Quasi-uniform meshes keep the reader's cell (the predicted gain is about two,
which buys nothing measurable on a step); structure-resolving meshes take a smaller one: a coastal
FESOM-C mesh goes from \ang{0.058} to \ang{0.014} (search 1.7 times faster) and an 11.5-million-node
global mesh from \ang{0.145} to \ang{0.051} (1.7--2.6 times). An explicit
\conf{mesh.grid\_hash\_cell\_size} always wins; the chosen cell and the rule are written to the
output metadata.

\paragraph{The node coordinate table.} The scans read unit-sphere coordinates of the nodes from a
table filled once when the hash is built, instead of evaluating trigonometric functions per
candidate. The pure scan is 6.3 times faster (\SI{3.2}{\micro s} to \SI{0.5}{\micro s} per query on one
thread), a two-dimensional midpoint step phase about twice as fast, at 24 bytes per node and process;
the trajectories are bit-for-bit unchanged.

\paragraph{The cost of exact anchors.} \Cref{tab:perf-guard} gives the cost of the guard for
\num{400000} particles at eight threads. It is a few per cent on HEALPix meshes and on the fused
three-dimensional kernel, about 40\,\% on a structured mid-latitude mesh even with every anchor provable,
and grows where many windows must widen --- at high latitude, where a degree of longitude shrinks, and with
a cell well below the node spacing. In every case the fused kernel with exact anchors stays faster
than the unfused path.

\begin{table}[t]
\centering\small
\caption{Cost of exact nearest-node anchors: one fused step, \num{400000} particles, eight threads,
milliseconds, with the node coordinate table. ``Off'' is the approximate in-window anchor; ``unfused''
is the unfused path under the guard.}
\label{tab:perf-guard}
\begin{tabular}{@{}lrrr@{}}
\toprule
Configuration & off & on & unfused \\
\midrule
structured mid-latitude mesh, automatic cell, RK4 & 115 & 157 & 230 \\
structured mesh at \ang{80}--\ang{88}\,N, automatic cell, RK4 & 102 & 201 & 278 \\
same, cell three times below the node spacing & 81 & 371 & 455 \\
HEALPix, polar band \ang{89}--\ang{89.9}\,N, RK4 barycentric & 531 & 855 & 961 \\
HEALPix, global uniform seeds, RK4 barycentric & 554 & 541 & 1019 \\
HEALPix at \ang{80}--\ang{88}\,N, fused 3-D midpoint barycentric & 398 & 385 & 406 \\
\bottomrule
\end{tabular}
\end{table}

\codeanchor{Cell rule: \fn{choose\_grid\_hash\_cell}, \fn{apply\_grid\_hash\_cell\_rule} and the
estimate \fn{\_gh\_sample\_stats}; table: \fn{\_node\_xyz\_fill}; exact search:
\fn{\_grid\_hash\_nearest\_exact}; all in \srcfile{core/mesh.py}. Benchmarks:
\srcfile{benchmarks/exact\_nearest\_cost.py}, \srcfile{benchmarks/nearest\_scan\_table.md},
\srcfile{benchmarks/grid\_hash\_cell\_rule\_table.md}.}

\section{Threads and processes on one node}\label{sec:perf-threads}

The kernels parallelize over particles. Two-dimensional fused kernels with single-precision
velocities scale steeply, about 38-fold for the midpoint method from one to 128 threads on a global
mesh. The fused three-dimensional midpoint kernel is not memory-bound: it gains 6.2-fold from 16 to
128 threads. The unfused three-dimensional paths are bound by memory bandwidth and saturate at 16 to
32 threads. The thread count must always be set explicitly (\code{NUMBA\_NUM\_THREADS}); left
unset, the kernel library occupies every logical core and small kernels slow down markedly.

A node can also be split into several processes, each tracking a contiguous share of the
particles and pinned to its own NUMA domain. \Cref{tab:perf-layouts} shows the effect for six million
particles in three dimensions on a zoom-9 HEALPix mesh with node-shared snapshots: the same 128 cores
take a 22\,\% shorter step as sixteen processes of eight threads than as one process of 128 (29\,\% more
particle-steps per second).

\begin{table}[t]
\centering\small
\caption{Process layouts on one node: 6.0 million particles, fused three-dimensional midpoint
barycentric kernel, zoom-9 HEALPix mesh (2.2 million wet surface nodes), 128 threads in total,
node-shared snapshots. Steady tracking step of the slowest process.}
\label{tab:perf-layouts}
\begin{tabular}{@{}lccccc@{}}
\toprule
Processes $\times$ threads & $1\times128$ & $2\times64$ & $4\times32$ & $8\times16$ & $16\times8$ \\
\midrule
Tracking step (s) & 0.576 & 0.540 & 0.489 & 0.478 & 0.447 \\
Throughput ($10^6$ particle-steps/s) & 10.4 & 11.1 & 12.3 & 12.6 & 13.4 \\
\bottomrule
\end{tabular}
\end{table}

\section{Mesh size and particle count}\label{sec:perf-mesh}

The kernel does not see the size of the mesh: a particle's work is a constant-size stencil around
its position. \Cref{tab:perf-mesh} confirms this for the fused three-dimensional step over a
3765-fold range of HEALPix meshes with the same six million particles --- the step varies by at most
1.45 times, and not monotonically; the search window holds a median of 57--65 nodes on all seven
meshes, so the variation is not the nearest-node search. The step grows slightly less than linearly
with the number of particles (\cref{fig:perf-particles}): from \SI{0.144}{\micro s} per particle-step at
three million to \SI{0.106}{\micro s} at 24 million on 128 threads, a log--log slope of 0.88.

\begin{table}[t]
\centering\small
\caption{Mesh size: 6.0 million particles, one 128-thread process, fused three-dimensional midpoint
barycentric kernel, HEALPix meshes of zoom 5 to 11.}
\label{tab:perf-mesh}
\begin{tabular}{@{}lrrrrrrr@{}}
\toprule
Zoom & 5 & 6 & 7 & 8 & 9 & 10 & 11 \\
\midrule
Wet surface nodes ($10^6$) & 0.009 & 0.036 & 0.14 & 0.57 & 2.2 & 8.9 & 35.5 \\
Tracking step (s) & 0.742 & 0.814 & 0.628 & 0.803 & 0.638 & 0.635 & 0.908 \\
\bottomrule
\end{tabular}
\end{table}

\begin{figure}[t]
\centering
\begin{tikzpicture}
\begin{loglogaxis}[width=0.62\linewidth,height=55mm,xlabel={particles},ylabel={tracking step (s)},
  grid=major,grid style={gray!20},tick label style={font=\footnotesize},label style={font=\small}]
\addplot[mark=*,color=lagblue,mark size=1.6pt] coordinates {
  (0.25e6,0.042) (0.5e6,0.100) (1.0e6,0.135) (1.5e6,0.253) (2.0e6,0.329) (3.0e6,0.433)
  (4.0e6,0.543) (6.0e6,0.794) (12e6,1.295) (24e6,2.546)};
\addplot[dashed,color=gray,domain=0.25e6:24e6,samples=2] {0.106e-6*x*(24e6/x)^(0.12)};
\end{loglogaxis}
\end{tikzpicture}
\caption{Particle count: steady tracking step against the number of particles, zoom-10 HEALPix mesh
(8.9 million wet surface nodes), one 128-thread process, fused three-dimensional midpoint kernel.
Dashed: the fitted power law with slope 0.88.}
\label{fig:perf-particles}
\end{figure}

\section{Many nodes}\label{sec:perf-nodes}

Particles are independent, so a release splits over nodes as contiguous shares, one or several
processes per node, each reading its own snapshots or sharing them within the node
(\cref{sec:perf-io}). \Cref{fig:perf-nodes} shows weak and strong scaling on the zoom-11 HEALPix mesh
(35.5 million wet surface nodes, 69 layers), one 128-thread process per node --- a zoom-11 process
needs about \SI{146}{GiB}, so one fits a \SI{256}{GB} node. With six million particles per node the
step stays within 90\,\% of the one-node value up to sixteen nodes (96 million particles in
\SI{0.85}{s}, $113\times10^{6}$ particle-steps per second). A fixed release of 24 million particles
runs 12.6 times faster on sixteen nodes; on the zoom-9 mesh, split as eight processes per node, six
million particles run 9.7 times faster on sixteen nodes, where each process holds only
\num{47000} particles.

Very large releases keep the same throughput. Tracked for a few steps without any output, $10^{9}$
particles on sixteen nodes step in \SI{7.6}{s} and $10^{10}$ particles on sixteen \SI{512}{GB} nodes in
\SI{73}{s} --- both about $1.3\times10^{8}$ particle-steps per second.

\begin{figure}[t]
\centering
\begin{tikzpicture}
\begin{axis}[name=weak,width=0.47\linewidth,height=52mm,xmode=log,log basis x=2,
  xtick={1,2,4,8,16},xticklabels={1,2,4,8,16},xlabel={nodes},ylabel={tracking step (s)},
  ymin=0,ymax=1.2,grid=major,grid style={gray!20},title={weak: 6.0\,M per node},
  title style={font=\small},tick label style={font=\footnotesize},label style={font=\small}]
\addplot[mark=*,color=lagblue,mark size=1.6pt] coordinates {(1,0.760) (2,0.792) (4,0.649) (8,0.829) (16,0.848)};
\addplot[dashed,color=gray,domain=1:16,samples=2] {0.760};
\end{axis}
\begin{axis}[at={(weak.east)},anchor=west,xshift=14mm,width=0.47\linewidth,height=52mm,
  xmode=log,ymode=log,log basis x=2,log basis y=2,xtick={1,2,4,8,16},xticklabels={1,2,4,8,16},
  ytick={1,2,4,8,16},yticklabels={1,2,4,8,16},xlabel={nodes},ylabel={speed-up},
  grid=major,grid style={gray!20},title={strong: 24.1\,M in total},title style={font=\small},
  tick label style={font=\footnotesize},label style={font=\small}]
\addplot[mark=*,color=lagblue,mark size=1.6pt] coordinates {(1,1.00) (2,1.99) (4,3.42) (8,9.02) (16,12.55)};
\addplot[dashed,color=gray,domain=1:16,samples=2] {x};
\end{axis}
\end{tikzpicture}
\caption{Scaling over nodes on the zoom-11 HEALPix mesh, one 128-thread process per node, fused
three-dimensional midpoint kernel. Left: weak scaling (dashed, the one-node step). Right: strong
scaling of 24.1 million particles (dashed, ideal). The strong curve is not smooth because the cost of
a process depends on which particles its contiguous share holds.}
\label{fig:perf-nodes}
\end{figure}

\section{Memory}\label{sec:perf-memory}

A process holds the mesh, the node coordinate table (24 bytes per node), the hash offsets (at most
\SI{128}{MB} under the cell rule), the snapshots of the read window, and the particle state. For
releases large enough that the particles dominate, the peak resident memory of a
three-dimensional midpoint process without output was measured as
\begin{equation}\label{eq:perf-memory}
  M_{\text{peak}} \approx \SI{61}{GiB} + \SI{181}{B} \times N_{\text{process}}
  \qquad\text{(zoom-11 mesh, fused 3-D midpoint, no output)},
\end{equation}
fitted through one-node releases of $10^{9}$ (\SI{230}{GiB}) and $4\times10^{9}$ particles
(\SI{735}{GiB}); it predicted \SI{166.7}{GiB} for 625 million particles per process where
\SI{166.2}{GiB} was measured. Below about $4\times10^{8}$ particles per process the snapshot read,
not the particles, sets the peak (122--146 GiB on this mesh). A small probe therefore measures the
read, not the cost per particle; a new configuration is sized from two releases large enough for the
particles to dominate.

\section{Input and output}\label{sec:perf-io}

\paragraph{The read term dominates runs that write.} For a full model year of six million
three-dimensional particles on a zoom-9 mesh with daily snapshots and six-hour steps, one 128-thread
process took 2 hours 12 minutes: of the mean step of \SI{5.33}{s}, reading took \SI{2.94}{s},
tracking \SI{0.62}{s} and writing \SI{1.78}{s} (\SI{7.1}{s} per record, every fourth step). Input and
output make up 88\,\% of the step. On the zoom-11 mesh one daily snapshot costs about \SI{146}{s} per
process --- as much as some 190 tracking steps of six million particles.

\paragraph{Levers on the read term.} Confining a study to part of the ocean does not by itself
reduce the read on a global mesh; extracting the region into a smaller mesh and forcing once, before
the run, reduces the input volume by about two orders of magnitude and improves the thread scaling,
because the smaller field fits in cache. The read-ahead (\cref{sec:temp-prefetch}) amortizes the
per-read overhead; it is forced to one snapshot when snapshots are shared between processes or when
the forcing is cycled. Several processes on one node can share one copy of each processed snapshot
(\conf{numerics.shared\_fields\_dir}): one process, elected by a file lock, publishes each snapshot
with an atomic rename and the others map it read-only. Every measurement placed the cache on the
parallel file system: shared within one node it was nearly free and removed the per-process replication
(\SI{16}{GB} to \SI{10}{GB} per process in one layout); shared by several nodes it slowed every step,
because remote nodes fault pages in while stepping. Node-local memory has not been measured. The outputs are bit-for-bit identical with and without the shared
cache, and any failure of the shared layer falls back to local reads.

\paragraph{Output.} The volume grows with the number of particles and records: at most about \SI{66}{B} per
particle and record uncompressed (positions, status, velocity diagnostics and two sampled fields, the
velocity diagnostics some 40\,\%), and about \SI{37}{B} on disk in the one measured compressed store
(six million particles, one year); \conf{output\_freq} is the
lever, and sampled variables add their own columns. A pulsed release stores the whole pre-allocated
array at every record, released or not.

\codeanchor{Read window and shared snapshots: \fn{\_VelCache} in \srcfile{tracking.py},
\fn{SharedFieldCache} in \srcfile{core/shared\_fields.py}; region extraction:
\srcfile{utils/extract\_fesom\_region.py} and \srcfile{utils/extract\_healpix\_region.py}; sharding:
\srcfile{utils/sort\_seeds.py}, \srcfile{utils/split\_seeds.py}, \srcfile{utils/concat\_trajectories.py}.}

\section{Sharding a release}\label{sec:perf-shard}

A release is sharded in three steps: the seeds are sorted spatially (by HEALPix cell, then depth),
so that consecutive particles share cached velocity data --- worth 17\,\% on the fused
three-dimensional step at a million particles; the sorted file is split into contiguous shares, one
per process; and the per-process outputs are merged in shard order afterwards. A deterministic run
with a single release, sharded and merged, is data-identical to the unsharded run, stochastic features
included: the splitting tool writes each shard's position in the full release into the shard file, and
the engine keys every random draw by the particle's index in the full release (\cref{sec:diff-rng}).
All shards therefore use the same seeds.

\section{Measuring the cost}\label{sec:perf-measure}

Three measurements cover most questions. The verbose per-step phase split of
\cref{sec:perf-model} tells which term dominates a run; the steady step --- the median tracking
phase over read-free steps after the first, of the slowest process --- is the number to compare
between runs and configurations, never a mean that mixes steps that read snapshots with steps that
do not. The portable single-node benchmark (\srcfile{benchmarks/single\_node/run\_bench.py}) builds a
synthetic HEALPix store and reports the machine-independent \emph{thread-microseconds per
particle}, its steady step (the mean of its last five steps) times the thread count divided by the
particles: on the compute node
above 15.7--19.9 thread-\si{\micro s} from 16 to 128 threads, on a laptop's five performance cores
5.5. Comparisons are valid only between rows of the same benchmark.

\codeanchor{Phase timing: the verbose branch of \fn{run\_single\_realization} in
\srcfile{tracking.py}; benchmark: \srcfile{benchmarks/single\_node/run\_bench.py} and its
\srcfile{README.md}; the diagnosis path for a run slower than these numbers is the
\code{kelp-diagnose} skill (\cref{ch:skills}).}
